% TiMBL 7.0 API

\documentclass{report}
\usepackage{a4wide}
\usepackage{palatino}
\usepackage{url}

% allow long URLs to break across lines
\urldef{\apimanualurl}\url{https://github.com/LanguageMachines/timbl/blob/master/docs/Timbl_7.0_API.pdf}
\urldef{\timblmanualurl}\url{https://github.com/LanguageMachines/timbl/blob/master/docs/Timbl_7.0_Manual.pdf}

\newcommand{\chisq}{{$ \chi^2 $}}

\author{Ko van der Sloot\\ \ \\
        Centre for Language Studies \\
        Radboud University Nijmegen \\ \ \\
        URL: https://languagemachines.github.io/timbl\thanks{This
        document is available from \apimanualurl.}}

\title{{\huge TiMBL: Tilburg Memory-Based Learner} \\ \vspace*{0.5cm}
{\bf version 7.0} \\ \vspace*{0.5cm}{\huge API Reference Guide}}

%better paragraph indentation
\parindent 0pt
\parskip 9pt


\begin{document}

\maketitle

\tableofcontents

\chapter*{Preface}

This is a brief description of the TimblAPI class, the application
programming interface to the
Timbl\footnote{\url{https://languagemachines.github.io/timbl}}
software package, and its main functions. For an introduction into
Timbl, consult the Timbl Reference Guide \cite{Daelemans+26}. Although
most of the API can be traced in the {\tt TimblAPI.h} file, the
reverse is not true; some functions in {\tt TimblAPI.h} are still
``work in progress'' and some others are artefacts to simplify the
implementation of the TiMBL main program\footnote{Timbl.cxx is
  therefore {\em not} a good example of how to use the API.}.

To learn more about using the API, you should study programs such as
{\tt classify.cxx}, {\tt tse.cxx}, and the examples given in this
manual, which can all be found in the {\tt demos} directory of this
distribution. As you can readily gather from these examples, the basic
thing you need to do to get access to the TimblAPI functions is to
include {\tt timbl/TimblAPI.h} in the program, and to link against the
{\tt libtimbl} library (the command {\tt pkg-config --cflags --libs
  timbl} gives the appropriate compiler and linker flags).

{\bf Important note}: The described functions return a result (mostly
a bool) to indicate succes or failure. To simplify the examples, we
ignore these return values. This is, of course, bad practice, to be avoided in
real life programming.\footnote{as stated by commandment 6 of ``The
  Ten Commandments for C Programmers''' by Henry Spencer:

If a function be advertised to return an error code in the event of
difficulties, thou shalt check for that code, yea, even though the
checks triple the size of thy code and produce aches in thy typing
fingers, for if thou thinkest ``it cannot happen to me'', the gods
shall surely punish thee for thy arrogance.}

{\bf Warning}: Although the TiMBL internals perform some sanity
checking, it is quite possible to combine API functions such
that some undetermined state is reached, or even a conflict
arises. The effect of the {\tt SetOptions()} function, for instance,
might be quite surprising. If you have created your own program
with the API it might be wise to test against well-know data to see if
the results make sense.

\chapter{Changes}
\label{changes}

\section{From version 6.3 to 7.0}

Between versions 6.3 and 7.0, TiMBL saw a long series of incremental
releases. Version 6.8 brought a major refactoring of the code base
that breaks both the API and the ABI; code written against earlier
versions may need to be adapted, and must in any case be
recompiled. Building against TiMBL now requires a C++17 compiler. The
changes that affect API users are:

\begin{itemize}
\item TiMBL is now Unicode-based throughout: input is interpreted as
  UTF-8 and internally normalised to NFC. Most classification
  functions now exist in two variants: one taking UTF-8 encoded {\tt
    std::string} arguments (as before), and one taking {\tt
    icu::UnicodeString} arguments. The {\tt Increment()} and {\tt
    Decrement()} functions have {\tt icu::UnicodeString} counterparts
  named {\tt Increment\_u()} and {\tt Decrement\_u()}.
\item The {\tt ValueDistribution} and {\tt WValueDistribution}
  classes have been renamed to {\tt ClassDistribution} and {\tt
    WClassDistribution}. Type aliases with the old names are provided
  for backward compatibility. The distribution is now stored as a
  flat sorted vector instead of a map, which changes the iterator
  interface: a {\tt dist\_iterator} now dereferences to a {\tt
    Vfield} directly, so you write {\tt it->Value()} and {\tt
    it->Weight()} where you used to write {\tt it->second->Value()}
  and {\tt it->second->Weight()}. See Section~\ref{advanced}.
\item The {\tt StartServer()} function has been removed. All server
  functionality lives in the separate TimblServer
  package\footnote{\url{https://github.com/LanguageMachines/timblserver}},
  which wraps trained TiMBL experiments in a TCP, HTTP or JSON
  server. See Chapter~\ref{server}.
\item The {\tt ShowBestNeighbors()} function lost its second (bool)
  parameter; whether distributions are shown now follows from the
  verbosity settings.
\item {\tt CurrentWeightings()} was renamed to {\tt
    GetCurrentWeights()}.
\item New functions were added, reflecting new functionality of
  Timbl 6.4--7.0:
  \begin{itemize}
  \item {\tt Prune()} prunes a trained IB1-like instance base to an
    IGTree, the API equivalent of the new {\tt --prune} command line
    option (see Section~\ref{prune});
  \item {\tt SetThreads()} sets the number of threads used for
    testing, the equivalent of {\tt --clones};
  \item {\tt confidence()} returns the confidence of the last
    classification (the equivalent of the {\tt +vcf} verbosity
    option, in combination with {\tt -G});
  \item {\tt NS\_Test()} bulk-tests a file using neighborSets;
  \item {\tt WriteInstanceBaseLevels()} dumps only the top levels of
    the instance base;
  \item {\tt Expand()} and {\tt Remove()} add or remove instances
    from a file to/from a trained instance base;
  \item {\tt lastHandledInstance()}, {\tt myTargets()}, {\tt
    getInputFormat()}, {\tt isValid()}, {\tt initExperiment()} and
    {\tt BuildInfo()} give access to more of TiMBL's internal state;
  \item a constructor taking a pre-parsed {\tt TiCC::CL\_Options}
    object was added next to the string-based constructor.
  \end{itemize}
\item The {\tt Weighting} enum gained the value {\tt SD} (standard
  deviation weighting for numeric features, the {\tt -w SD} option).
\end{itemize}

\section{From version 6.2 to 6.3}

No changes to the API are made for this release. This Manual is made
up to date (preserving the beta-state).

\section{From version 6.1 to 6.2}

In version 6.2, some additional functions were added to the API: {\tt
  matchDepth()}, {\tt matchedAtLeaf()}, {\tt WriteMatrices()}, {\tt
  GetMatrices()} and {\tt ShowStatistics()}. These reflect the
additional functionality of Timbl 6.2.  The API is still experimental,
and contains more functions than described in this manual. Using these
`undocumented' features is, as usual, unwise.

\section{From version 5.1 to 6.1}

The major change in 6.0 is the introduction of the {\tt neighborSet}
class, with some special Classify functions.  We added Classify
functions that deliver pointers into Timbl's internal data. This is
fast, but dangerous.  Also, a {\tt WriteInstanceBaseXml()} function is
added, which comes in handy when you want to know more about the
instance base.  Two more examples demonstrating neighborSets and such
are added in Appendix B. From version 6.0 to 6.1, the API has not changed.

\section{From version 5.0 to 5.1}

The API is quite stable at the moment. Most TiMBL changes did not
affect the API. The only real API change is in the {\tt GetWeights()}
function. (see the section on Storing and retrieving intermediate
results).

\chapter{Quick-start}
\section{Setting up an experiment}

There is just one way to start a TiMBL experiment, which is to call
the TimblAPI constructor:

\begin{footnotesize}
\begin{verbatim}
  TimblAPI( const std::string& args, const std::string& name ="" );
\end{verbatim}
\end{footnotesize}

args is used as a "command line" and is parsed for all kind of options
which are used to create the right kind of experiment with the desired
settings for metric, weighting etc. If something is wrong with the
settings, {\em no}\/ object is created.

The most important option is {\tt -a}  to set the kind of algorithm,
e.g. {\tt -a IB1} to invoke an IB1 experiment or {\tt -a IGTREE} to invoke an IGTREE
experiment. All command line options are described in the TiMBL
Reference Guide \cite{Daelemans+26}.

The optional name can be useful if you have multiple experiments.
In case of warnings or errors, this name is appended to the message.

For example:

\begin{footnotesize}
\begin{verbatim}
  TimblAPI *My_Experiment = new TimblAPI( "-a IGTREE +vDI+DB",
                                          "test1" );
\end{verbatim}
\end{footnotesize}

{\tt My\_Experiment} is created as an IGTREE experiment with the name
"test1", and the verbosity is set to DI+DB, meaning that the output
will contain DIstance and DistriBution information.

A second constructor takes a pre-parsed {\tt TiCC::CL\_Options}
object instead of an options string; this is convenient when your
program passes (a subset of) its own command line on to TiMBL.

The counterpart to creation is the {\tt \~{ }TimblAPI()} destructor,
which is called when you delete an experiment:

\begin{footnotesize}
\begin{verbatim}
  delete My_Experiment;
\end{verbatim}
\end{footnotesize}

\section{Running an experiment}

Assuming that we have appropriate datafiles (such as the example files {\tt
dimin.train} and {\tt dimin.test} in the TiMBL package), we can get
started right away with the functions {\tt Learn()} and {\tt Test()}.

\subsection{Training}
\begin{footnotesize}
\begin{verbatim}
  bool Learn( const std::string& f );
\end{verbatim}
\end{footnotesize}

This function takes a file with name 'f', and gathers information
such as: number of features, number and frequency of feature values and
the same for class names. After that, these data are used to calculate
a lot of statistical information, which will be used for
testing. Finally, an InstanceBase is created, tuned to the current
algorithm.

\subsection{Testing}
\begin{footnotesize}
\begin{verbatim}
  bool Test( const std::string& in,
             const std::string& out,
             const std::string& perc = "" );
\end{verbatim}
\end{footnotesize}

Test a file given by 'in' and write results to 'out'. If 'perc' is not
empty, then a percentage score is written to file 'perc'.

For example:

\begin{footnotesize}
\begin{verbatim}
  My_Experiment->Learn( "dimin.train" );
  My_Experiment->Test( "dimin.test", "my_first_test" );
\end{verbatim}
\end{footnotesize}

An InstanceBase will be created from dimin.train, then dimin.test is
tested against that InstanceBase and output is written to
my\_first\_test.

Testing can be parallelised over several threads with

\begin{footnotesize}
\begin{verbatim}
  bool SetThreads( int c );
\end{verbatim}
\end{footnotesize}

which is the API equivalent of the {\tt --clones} command line option.

\subsection{Special cases of {\tt Learn()} and {\tt Test()}}

There are special cases where {\tt Learn()} behaves differently:

\begin{itemize}
\item When the algorithm is IB2, {\tt Learn()} will automatically take
  the first $n$ lines of f (set with the {\tt -b n} option) to
  bootstrap itself, and then the rest of f for IB2-learning. After
  Learning IB2, you can use {\tt Test()} as usual.

\item When the algorithm is CV, {\tt Learn()} is not defined, and all
  work is done in a special version of {\tt Test()}. 'f' is assumed to
  give the name of a file, which, on separate lines, gives the names
  of the files to be cross-validated.

  Also, if {\em featureWeights}\/ or {\em probabilities}\/ are read from
  user-defined datafiles, a special {\tt CVprepare()} function must be called,
  to make the weigthing, weightFilename and probabilityFileName known to the
{\tt Test()} function.

See Appendix B for a complete CV example (program {\tt api\_test3}).

\end{itemize}

\section{More about settings}

After an experiment is set up with the TimblAPI constructor, many
options can be changed "on the fly" with:

\begin{footnotesize}
\begin{verbatim}
  bool SetOptions( const std::string& opts );
\end{verbatim}
\end{footnotesize}

Here, `opts' is interpreted as a list of option settings, just like in
the TimblAPI constructor. When an error in the opts string is found,
{\tt SetOptions()} returns false. Whether any options are really set
or changed in that case is undefined. Note that a few options can only
be set {\em once}\/ when creating the experiment, most notably the
algorithm. Any attempt to change these options will result in a
failure.

Note: {\tt SetOptions()} is lazy; changes are cached until the
moment they are really needed, so you can do several {\tt SetOptions()}
calls with even different values for the same option. Only the last
one seen will be used for running the experiment.

To see which options are in effect, you can use the calls {\tt ShowOptions()}
and {\tt ShowSettings()}.

\begin{footnotesize}
\begin{verbatim}
  bool ShowOptions( std::ostream& );
\end{verbatim}
\end{footnotesize}

Shows all options with their possible and current values.

\begin{footnotesize}
\begin{verbatim}
  bool ShowSettings( std::ostream& );
\end{verbatim}
\end{footnotesize}

Shows all options and their currect values.

For example:

\begin{footnotesize}
\begin{verbatim}
  My_Experiment->SetOptions( "-w2 -m:M" );
  My_Experiment->SetOptions( "-w3 -v:DB" );
  My_Experiment->ShowSettings( cout )
\end{verbatim}
\end{footnotesize}

See Appendix B (program {\tt api\_test1}) for the output.

\section{Storing and retrieving intermediate results}

To speed up testing, or to manipulate what is happening internally, we
can store and retrieve several important parts of our experiment: The
InstanceBase, the FeatureWeights, the ProbabilityArrays and the ValueDistance Matrices.

Saving is done with:

\begin{footnotesize}
\begin{verbatim}
  bool WriteInstanceBase( const std::string& f );
  bool SaveWeights( const std::string& f );
  bool WriteArrays( const std::string& f );
  bool WriteMatrices( const std::string& f );
\end{verbatim}
\end{footnotesize}

Retrieve with their counterparts:

\begin{footnotesize}
\begin{verbatim}
  bool GetInstanceBase( const std::string& f );
  bool GetWeights( const std::string& f, Weighting w );
  bool GetArrays( const std::string& f );
  bool GetMatrices( const std::string& f );
\end{verbatim}
\end{footnotesize}

All use `f' as a filename for storing/retrieving. {\tt GetWeights} needs
information to decide {\em which}\/ weighting to retrieve.
Weighting is defined as the enumerated type:

\begin{footnotesize}
\begin{verbatim}
  enum Weighting { UNKNOWN_W, UD, NW, GR, IG, X2, SV, SD };
\end{verbatim}
\end{footnotesize}

Some notes:

\begin{enumerate}
\item The InstanceBase is stored in a internal format, with or without
hashing, depending on the {\tt -H} option. The format is described in the
TiMBL manual. Remember that it is a bad idea to edit this file in any way.
\item {\tt GetWeights()} can be used to override the weights that
{\tt Learn()} calculated. {\tt UNKNOWN\_W} should not be used.
\item The Probability arrays are described in the TiMBL manual. They can be
manipulated to tune the MVDM similarity metric.
\end{enumerate}

If you like you may dump the Instancebase in an XML format. No Retrieve
function is available for this format.

\begin{footnotesize}
\begin{verbatim}
  bool WriteInstanceBaseXml( const std::string& f );
\end{verbatim}
\end{footnotesize}

Related is:

\begin{footnotesize}
\begin{verbatim}
  bool WriteInstanceBaseLevels( const std::string& f, unsigned int l );
\end{verbatim}
\end{footnotesize}

which dumps only the first 'l' levels of the InstanceBase, which can
be useful to inspect the global structure of a large tree.

\chapter{Classify functions}

\section{Classify functions: Elementary}
After an experiment is trained with {\tt Learn()}, we do not have to use
{\tt Test()} to do bulk-testing on a file.
We can create our own tests with the {\tt Classify} functions:

\begin{footnotesize}
\begin{verbatim}
  bool Classify( const std::string& Line, std::string& result );
  bool Classify( const std::string& Line, std::string& result,
                 double& distance );
  bool Classify( const std::string& Line, std::string& result,
                 std::string& Distrib, double& distance );
\end{verbatim}
\end{footnotesize}

Results are stored in 'result' (the assigned class). 'distance' will
get the calculated distance, and 'Distrib' the distribution at
'distance' which is used to calculate 'result'.  Distrib will be a
string like ``\{ NP 2, PP 6 \}''. It is up to you to parse and
interpret this. (In this case: There were 8 classes assigned at
'distance', 2 NP's and 6 PP's, giving a 'result' of ``PP''.)

The {\tt std::string} arguments are interpreted as UTF-8 encoded
Unicode. Since version 6.8, a variant using ICU Unicode strings is
also available:

\begin{footnotesize}
\begin{verbatim}
  bool Classify( const icu::UnicodeString& Line,
                 icu::UnicodeString& result );
\end{verbatim}
\end{footnotesize}

If you want to perform analyses on these distributions, it might be a
good idea to read the next section about the other range of Classify()
functions.

A main disadvantage compared to using {\tt Test()} is that {\tt
  Test()} is optimized.  {\tt Classify()} has to test for sanity of
its input and also whether a {\tt SetOptions()} has been
performed. This slows down the process.

A good example of the use of {\tt Classify()} is the {\tt
 classify.cxx} program in the TiMBL Distribution.

Depending on the Algorithm and Verbosity setting, it may be possible
to get some extra information on the details of each classification
using:

\begin{footnotesize}
\begin{verbatim}
   bool ShowBestNeighbors( std::ostream& os ) const;
\end{verbatim}
\end{footnotesize}

Provided that the option {\tt +v n} or {\tt +v k} is set and we use
IB1 or IB2, output is produced similar to what we see in the TiMBL
program.  Bear in mind: The {\tt +vn} option is expensive in time
and memory and does not work for IGTREE, TRIBL, and TRIBL2.

Two other functions provide the results as given by the {\tt +vmd} verbosity
option:

\begin{footnotesize}
\begin{verbatim}
    size_t matchDepth() const;
    bool matchedAtLeaf() const;
\end{verbatim}
\end{footnotesize}

The first returns the matching Depth in the InstanceBase; the second
flags whether it was a Leaf or a Non-Terminal Node.

Finally, when a normalisation is set with the {\tt -G} option,

\begin{footnotesize}
\begin{verbatim}
    double confidence() const;
\end{verbatim}
\end{footnotesize}

returns the confidence of the last classification, as given by the
{\tt +vcf} verbosity option.

\section{Classify functions: Advanced}
\label{advanced}

A faster, but more dangerous version of Classify is also available.
It is faster because it returns pointers into Timbl's internal
datastructures. It is dangerous because it returns pointers into
Timbl's internal datastructures (using 'const' pointers, so it is
fortunately difficult te really damage Timbl)

\begin{footnotesize}
\begin{verbatim}
  const TargetValue *Classify( const std::string& );
  const TargetValue *Classify( const std::string&,
                               const ClassDistribution *& );
  const TargetValue *Classify( const std::string&, double& );
  const TargetValue *Classify( const std::string&,
                               const ClassDistribution *&,
                               double& );
\end{verbatim}
\end{footnotesize}

All four also exist in a variant taking a {\tt const
  icu::UnicodeString\&} as their first argument.

A ClassDistribution (called ValueDistribution before version 6.8; an
alias with that name is still provided) is a list-like object (but it
is not a real list!) that contains TargetValues objects and
weights. It is the result of combining all nearest neighbors and
applying the desired weightings.  Timbl chooses a best TargetValue
from this ClassDistribution and the Classify functions return that as
their main result.

{\bf Important}: Because these functions return pointers into Timbl's
internal representation, the results are only valid until the next
Classify function is called (or the experiment is deleted).

Both the TargetValue and ClassDistribution objects have output
operators defined, so you can print them.  TargetValue also has a {\tt
  name\_string()} function, which returns a std::string so you can
collect results.  ClassDistribution has an iterator-like interface
which makes it possible to walk through the Distribution.

An iterator on a {\tt ClassDistribution *vd} is created like this:
\begin{footnotesize}
\begin{verbatim}
  ClassDistribution::dist_iterator it=vd->begin();
\end{verbatim}
\end{footnotesize}

The iterator walks through the entries of the distribution. Every
entry is a {\tt Vfield} object; you may print it directly, or extract
its {\tt Value()} (which happens to be a TargetValue pointer) or
extract its {\tt Weight()}, which is a {\tt double}.

Like this:
\begin{footnotesize}
\begin{verbatim}
  while ( it != vd->end() ){
    cout << *it << " has a value: ";
    cout << it->Value() << " and a weight of "
         << it->Weight() << endl;
    ++it;
  }
\end{verbatim}
\end{footnotesize}

Printing {\tt *it} is the same as printing the
TargetValue plus its Weight. (Before version 6.8, the iterator
dereferenced to a pair, and you had to write {\tt it->second->Value()}
and {\tt it->second->Weight()}; the {\tt ->second} part is now gone.)

In the {\em demos}\/ directory you will find a complete example in api\_test6.

{\bf Warning}: it is possible to search the Timbl code for the
internal representation of the TargetValue and ClassDistribution
objects, but please DON'T DO THAT.  The representation might change
between Timbl versions.

\section{Classify functions: neighborSets}

A more flexible way of classifying is to use one of these functions:

\begin{footnotesize}
\begin{verbatim}
  const neighborSet *classifyNS( const icu::UnicodeString& );
  bool classifyNS( const icu::UnicodeString&, neighborSet& );
  bool classifyNS( const std::string&, neighborSet& );
\end{verbatim}
\end{footnotesize}

The first function will classify an instance and return a pointer to a
{\tt neighborSet} object. This object may be seen as an container
which holds both distances and distributions up to a certain depth,
(which is {\em at least}\/ the number of neighbors (-k option) that
was used for the classifying task.)  It is a const object, so you
cannot directly manipulate its internals, but there are some
functions defined to get useful information out of the neighborSet.

Important:  The neighborSet {\em will be overwritten}\/ on the next
call to any of the classify functions. Be sure to get all the
results out before that happens.

To make life easy, the other variants can be used, which fill a
neighborSet object that you provide (the same could be achieved by a
copy of the result of the first function); the {\tt std::string}
variant interprets its input as UTF-8.

{\bf Note}: NeighborSets can be large, and copying therefore
expensive, so you should only do this if you really have to.

There is also a neighborSet-based counterpart of {\tt Test()}:

\begin{footnotesize}
\begin{verbatim}
  bool NS_Test( const std::string& in, const std::string& out );
\end{verbatim}
\end{footnotesize}

which bulk-tests file 'in', writing per-instance neighborSet
information to file 'out'.

\subsection{How to get results from a neighborSet}

No metric functions (such as exponential decay and the like) are
performed on the neighborSet. You are free to insert your own metrics, or
use Timbls built-in metrics.

\begin{footnotesize}
\begin{verbatim}
  double getDistance( size_t n ) const;
  double bestDistance() const;
  const ClassDistribution *getDistribution( size_t n ) const;
  WClassDistribution *bestDistribution( const decayStruct * ds=0,
                                        size_t n=0 ) const ;
\end{verbatim}
\end{footnotesize}

{\tt getDistance( n )} will return the distance of the neighbor(s) at n.
{\tt bestDistance()} is simply {\tt getDistance(0)}.

{\tt getDistribution( n )} will return the distribution of neighbor(s) at
n.

{\tt bestDistribution()} will return the Weighted distribution
calculated using the first n elements in the container and a metric
specified by the {\tt decayStruct}.  The default n=0, means: use the
whole container. An empty decay struct means zeroDecay.

The returned WClassDistribution object is handed to you, and you are
responsible for deleting it after using it (see the previous section
for more details about ClassDistributions).

A decayStruct is one of:

\begin{footnotesize}
\begin{verbatim}
  class zeroDecay();
  class invLinDecay();
  class invDistDecay();
  class expDecay( double alpha );
  class expDecay( double alpha, double beta );
\end{verbatim}
\end{footnotesize}

For example, to get a WClassDistribution form a neighborSet {\tt nb}, using
3 neighbors and exponential decay with alpha=0.3, you can do:

\begin{footnotesize}
\begin{verbatim}
  decayStruct *dc = new  expDecay(0.3);
  WClassDistribution *vd = nb->bestDistribution( dc, 3 );
\end{verbatim}
\end{footnotesize}


\subsection{Useful operations on neighborSet objects}

You can print neighborSet objects:

\begin{footnotesize}
\begin{verbatim}
    std::ostream& operator<<( std::ostream&, const neighborSet& );
    std::ostream& operator<<( std::ostream&, const neighborSet * );
\end{verbatim}
\end{footnotesize}

You may create a neighborSet yourself, and assign and delete them:

\begin{footnotesize}
\begin{verbatim}
    neighborSet();
    neighborSet( const neighborSet& );
    neighborSet& operator=( const neighborSet& );
    ~neighborSet();
\end{verbatim}
\end{footnotesize}

If you create an neighborSet, you might want to reserve space for it,
to avoid needless reallocations. Also it can be cleared, and you can
ask the size (just like with normal containers):

\begin{footnotesize}
\begin{verbatim}
    void reserve( size_t );
    void clear();
    size_t size() const;
\end{verbatim}
\end{footnotesize}

Two neighborSets can be merged:

\begin{footnotesize}
\begin{verbatim}
    void merge( const neighborSet& );
\end{verbatim}
\end{footnotesize}

A neighborSet can be truncated at a certain level. This is useful
after merging neighborSets. Merging sets with depth k and n will
result in a set with a depth somewhere within the range $[max(k,n), k+n]$.

\begin{footnotesize}
\begin{verbatim}
    void truncate( size_t );
\end{verbatim}
\end{footnotesize}

\chapter{Advanced Functions}

\section{Modifying the InstanceBase}

The instanceBase can be modified with the functions:

\begin{footnotesize}
\begin{verbatim}
  bool Increment( const std::string& Line );
  bool Decrement( const std::string& Line );
  bool Increment_u( const icu::UnicodeString& Line );
  bool Decrement_u( const icu::UnicodeString& Line );
\end{verbatim}
\end{footnotesize}

These functions add an Instance (as described by Line) to the
InstanceBase, or remove it.  This can only be done for IB1-like
experiments (IB1, IB2, CV and LOO), and enforces a lot of
statistical recalculations. The {\tt \_u} variants take an ICU
Unicode string; the {\tt std::string} variants interpret their input
as UTF-8.

More sophisticated are:

\begin{footnotesize}
\begin{verbatim}
  bool Expand( const std::string& File  );
  bool Remove( const std::string& File );
\end{verbatim}
\end{footnotesize}

which use the contents of File to do a bulk of Increments or Decrements, and
recalculate afterwards.

\section{Pruning an InstanceBase to an IGTree}
\label{prune}

New in version 7.0 is the possibility to prune a trained (or read)
IB1 or TRIBL2 instance base into an IGTree, without retraining from
the original data:

\begin{footnotesize}
\begin{verbatim}
  bool Prune( bool restore_distributions = false );
\end{verbatim}
\end{footnotesize}

This is the API equivalent of the {\tt --prune} command line
option. A typical use is to read a stored instance base with {\tt
  GetInstanceBase()}, optionally {\tt Expand()} it with new material,
{\tt Prune()} it, and store the result with {\tt
  WriteInstanceBase()}. After pruning, the experiment behaves like an
IGTree experiment. See the TiMBL Reference Guide \cite{Daelemans+26}
for details on the pruning process.

\section{Getting more information out of Timbl}

There are a few convenience functions to get extra information on
TiMBL and its behaviour:

\begin{footnotesize}
\begin{verbatim}
  bool WriteNamesFile( const std::string& f );
\end{verbatim}
\end{footnotesize}

Create a file which resembles a C4.5 namesfile.

\begin{footnotesize}
\begin{verbatim}
  Algorithm Algo()
\end{verbatim}
\end{footnotesize}

Give the current algorithm as a type enum Algorithm. First, the
declaration of the Algorithm type:

\begin{footnotesize}
\begin{verbatim}
  enum Algorithm { UNKNOWN_ALG, IB1, IB2, IGTREE,
                   TRIBL, TRIBL2, LOO, CV };
\end{verbatim}
\end{footnotesize}

This can be printed with the helper function:

\begin{footnotesize}
\begin{verbatim}
  const std::string to_string( const Algorithm )
\end{verbatim}
\end{footnotesize}

\begin{footnotesize}
\begin{verbatim}
  Weighting CurrentWeighting()
\end{verbatim}
\end{footnotesize}

Gives the current weighting as a type enum Weighting.

Declaration of Weighting:

\begin{footnotesize}
\begin{verbatim}
  enum Weighting { UNKNOWN_W, UD, NW, GR, IG, X2, SV, SD };
\end{verbatim}
\end{footnotesize}

This can be printed with the helper function:

\begin{footnotesize}
\begin{verbatim}
  const std::string to_string( const Weighting )
\end{verbatim}
\end{footnotesize}


\begin{footnotesize}
\begin{verbatim}
  Weighting GetCurrentWeights( std::vector<double>& v )
\end{verbatim}
\end{footnotesize}

Returns the current weighting as a type enum Weighting and also a
vector v with all the current values of this weighting. (This
function was called {\tt CurrentWeightings()} before version 6.8.)

\begin{footnotesize}
\begin{verbatim}
  std::string ExpName()
\end{verbatim}
\end{footnotesize}

Returns the value of 'name' given at the construction of the experiment

\begin{footnotesize}
\begin{verbatim}
  static std::string VersionInfo( bool full = false )
\end{verbatim}
\end{footnotesize}

Returns a string containing the Version number, the Revision and the
Revision string of the current API implementation. If full is true,
also information about the date and time of compilation is included.

\chapter{Server mode}
\label{server}

Before version 6.4, the API contained a {\tt StartServer()}
function. All server functionality has been removed from TiMBL and
lives in the separate TimblServer
package\footnote{\url{https://github.com/LanguageMachines/timblserver}},
which wraps one or more trained TiMBL experiments in a network server
speaking a TCP, HTTP or JSON protocol. See the TimblServer manual for
details.

\clearpage
\chapter{Annotated example programs}

The example programs and the data files they expect ({\tt
  dimin.train}, {\tt dimin.test}, {\tt small\_?.train} and {\tt
  cross\_val.test}) can be found in the {\tt demos} directory of the
TiMBL source distribution.

\subsection{example 1, {\tt api\_test1.cxx}}
\begin{footnotesize}
\begin{verbatim}
#include "timbl/TimblAPI.h"
int main(){
  Timbl::TimblAPI My_Experiment( "-a IGTREE +vDI+DB+F", "test1" );
  My_Experiment.SetOptions( "-w3 -vDB" );
  My_Experiment.ShowSettings( std::cout );
  My_Experiment.Learn( "dimin.train" );
  My_Experiment.Test( "dimin.test", "my_first_test.out" );
  My_Experiment.SetOptions( "-mM" );
  My_Experiment.Test( "dimin.test", "my_first_test.out" );
}
\end{verbatim}
\end{footnotesize}


Output (note that the settings are listed alphabetically nowadays):
\begin{footnotesize}
\begin{verbatim}
Current Experiment Settings :
ALL_WEIGHTS          : false
BEAM_SIZE            : 0
BIN_SIZE             : 20
CLIP_FACTOR          : 10
DECAY                : Z
DECAYPARAM_A         : 1.00000
DECAYPARAM_B         : 1.00000
DO_DIVERSIFY         : false
DO_PRUNE             : false
DO_SILLY             : false
DO_SLOPPY_LOO        : false
EXACT_MATCH          : false
EXEMPLAR_WEIGHTS     : false
FLENGTH              : 0
GLOBAL_METRIC        : O
HANDLE_OCCURRENCES   : 0
HASHED_TREE          : true
IB2_OFFSET           : 0
IG_THRESHOLD         : 1000
IGNORE_EXEMPLAR_WEIGHTS : true
INPUTFORMAT          : Unknown
KEEP_DISTRIBUTIONS   : false
MAXBESTS             : 500
METRICS              :
MVD_LIMIT            : 1
NEIGHBORS            : 1
NO_EXEMPLAR_WEIGHTS_TEST : true
NORM_FACTOR          : 1.00000
NORMALISATION        : None
PROGRESS             : 100000
SEED                 : -1
TARGET_POS           : 18446744073709551615
TREE_ORDER           : Unknown
TRIBL_OFFSET         : 0
VERBOSITY            : F+DI                             [Note 1]
WEIGHTING            : x2                               [Note 2]

Examine datafile 'dimin.train' gave the following results:
Number of Features: 12
InputFormat       : C4.5

-test1-Phase 1: Reading Datafile: dimin.train
-test1-Start:          0 @ Fri Jul  3 10:00:22 2026
-test1-Finished:    2999 @ Fri Jul  3 10:00:22 2026
-test1-Calculating Entropy         Fri Jul  3 10:00:22 2026
Lines of data     : 2999
DB Entropy        : 1.6178929
Number of Classes : 5

Feats	Vals	X-square	Variance	InfoGain	GainRatio
1          3	128.41828	0.021410184	0.030971064	0.024891536
    2     50	364.75812	0.030406645	0.060860038	0.027552191
    3     19	212.29804	0.017697402	0.039562857	0.018676787
    4     37	449.83823	0.037499019	0.052541227	0.052620750
    5      3	288.87218	0.048161417	0.074523225	0.047699231
    6     61	415.64113	0.034648310	0.10604433	0.024471911
    7     20	501.33465	0.041791818	0.12348668	0.034953203
    8     69	367.66021	0.030648567	0.097198760	0.043983864
    9      2	169.36962	0.056475363	0.045752381	0.046816705
   10     64	914.61906	0.076243669	0.21388759	0.042844587
   11     18	2807.0418	0.23399815	0.66970458	0.18507018
   12     43	7160.3682	0.59689631	1.2780762	0.32537181

-test1-Preparation took 0 seconds, 7 milliseconds and 343 microseconds
Feature Permutation based on Chi-Squared :
< 12, 11, 10, 7, 4, 6, 8, 2, 5, 3, 9, 1 >
-test1-Phase 2: Building multi index on Datafile: dimin.train
-test1-Start:          0 @ Fri Jul  3 10:00:22 2026
-test1-Finished:    2999 @ Fri Jul  3 10:00:22 2026
-test1-
Phase 3: Learning from Datafile: dimin.train
-test1-Start:          0 @ Fri Jul  3 10:00:22 2026
-test1-Finished:    2999 @ Fri Jul  3 10:00:22 2026

Size of InstanceBase = 148 Nodes, (5920 bytes), 99.61 % compression
-test1-Learning took 0 seconds, 25 milliseconds and 428 microseconds
Examine datafile 'dimin.test' gave the following results:
Number of Features: 12
InputFormat       : C4.5


Starting to test, Testfile: dimin.test
Writing output in:          my_first_test.out
Algorithm     : IGTree
Weighting     : Chi-square
Feature 1	 : 128.418283576224439
Feature 2	 : 364.758115277811896
Feature 3	 : 212.298037236345095
Feature 4	 : 449.838231470681876
Feature 5	 : 288.872176256387263
Feature 6	 : 415.641126446691771
Feature 7	 : 501.334653478280984
Feature 8	 : 367.660212489714240
Feature 9	 : 169.369615106487458
Feature 10	 : 914.619058199288816
Feature 11	 : 2807.041753278295346
Feature 12	 : 7160.368151902808677

-test1-Tested:      1 @ Fri Jul  3 10:00:22 2026
-test1-Tested:      2 @ Fri Jul  3 10:00:22 2026
-test1-Tested:      3 @ Fri Jul  3 10:00:22 2026
-test1-Tested:      4 @ Fri Jul  3 10:00:22 2026
-test1-Tested:      5 @ Fri Jul  3 10:00:22 2026
-test1-Tested:      6 @ Fri Jul  3 10:00:22 2026
-test1-Tested:      7 @ Fri Jul  3 10:00:22 2026
-test1-Tested:      8 @ Fri Jul  3 10:00:22 2026
-test1-Tested:      9 @ Fri Jul  3 10:00:22 2026
-test1-Tested:     10 @ Fri Jul  3 10:00:22 2026
-test1-Tested:    100 @ Fri Jul  3 10:00:22 2026
-test1-Ready:     950 @ Fri Jul  3 10:00:22 2026
Seconds taken: 0.0719 (13214.63 p/s)

overall accuracy:        0.962105  (914/950)
Examine datafile 'dimin.test' gave the following results:
Number of Features: 12
InputFormat       : C4.5

Warning:-test1-Metric must be Overlap for IGTree test.     [Note 3]
\end{verbatim}
\end{footnotesize}


Notes:
\begin{enumerate}
\item The first {\tt SetOptions()} sets the verbosity with {\tt +F+DI+DB}.
The second {\tt SetOptions()}, however, sets the verbosity with {\tt -vDB}, and the resulting verbosity is therefore {\tt F+DI}.
\item The {\tt -w2} of the first {\tt SetOptions()} is overruled with
  {\tt -w3} from the second {\tt SetOptions()}, resulting in a
  weighting of 3 or Chi-Square.
\item Due to the second {\tt SetOptions()}, the default metric is set to
MVDM --- this is however not applicable to IGTREE. This raises a warning
when we start to test.
\end{enumerate}

Result in my\_first\_test.out (first 20 lines):
\begin{footnotesize}
\begin{verbatim}
=,=,=,=,=,=,=,=,+,p,e,=,T,T        6619.8512628162
=,=,=,=,+,k,u,=,-,bl,u,m,E,P        2396.8557978603
+,m,I,=,-,d,A,G,-,d,},t,J,J        6619.8512628162
-,t,@,=,-,l,|,=,-,G,@,n,T,T        6619.8512628162
-,=,I,n,-,str,y,=,+,m,E,nt,J,J        6619.8512628162
=,=,=,=,=,=,=,=,+,br,L,t,J,J        6619.8512628162
=,=,=,=,+,zw,A,=,-,m,@,r,T,T        6619.8512628162
=,=,=,=,-,f,u,=,+,dr,a,l,T,T        6619.8512628162
=,=,=,=,=,=,=,=,+,l,e,w,T,T        13780.219414719
=,=,=,=,+,tr,K,N,-,k,a,rt,J,J        6619.8512628162
=,=,=,=,+,=,o,=,-,p,u,=,T,T        3812.8095095379
=,=,=,=,=,=,=,=,+,l,A,m,E,E        3812.8095095379
=,=,=,=,=,=,=,=,+,l,A,p,J,J        6619.8512628162
=,=,=,=,=,=,=,=,+,sx,E,lm,P,P        6619.8512628162
+,l,a,=,-,d,@,=,-,k,A,st,J,J        6619.8512628162
-,s,i,=,-,f,E,r,-,st,O,k,J,J        6619.8512628162
=,=,=,=,=,=,=,=,+,sp,a,n,T,T        6619.8512628162
=,=,=,=,=,=,=,=,+,st,o,t,J,J        6619.8512628162
=,=,=,=,+,sp,a,r,-,b,u,k,J,J        6619.8512628162
+,h,I,N,-,k,@,l,-,bl,O,k,J,J        6619.8512628162
\end{verbatim}
\end{footnotesize}
\clearpage

\subsection{example 2, {\tt api\_test2.cxx}}

This demonstrates IB2 learning. Our example program:

\begin{footnotesize}
\begin{verbatim}
#include <cstdlib>
#include "timbl/TimblAPI.h"
int main(){
  Timbl::TimblAPI *My_Experiment
    = new Timbl::TimblAPI( "-a IB2 +vF+DI+DB" , "test2" );
  My_Experiment->SetOptions( "-b100" );
  My_Experiment->ShowSettings( std::cout );
  My_Experiment->Learn( "dimin.train" );
  My_Experiment->Test( "dimin.test", "my_second_test.out" );
  delete My_Experiment;
  exit(1);
}
\end{verbatim}
\end{footnotesize}

We create an experiment for the IB2 algorithm, with the {\tt -b} option set
to 100, so the first 100 lines of {\tt dimin.train} will be used to
bootstrap the learning, as we can see from the output (the initial
settings dump is omitted here; it is analogous to the one in example
1, with {\tt IB2\_OFFSET : 100}):

\begin{footnotesize}
\begin{verbatim}
Examine datafile 'dimin.train' gave the following results:
Number of Features: 12
InputFormat       : C4.5

-test2-Phase 1: Reading Datafile: dimin.train
-test2-Start:          0 @ Fri Jul  3 10:00:23 2026
-test2-Finished:    2999 @ Fri Jul  3 10:00:23 2026
-test2-Calculating Entropy         Fri Jul  3 10:00:23 2026
Lines of data     : 2999                                  [Note 1]
DB Entropy        : 1.6178929
Number of Classes : 5

Feats	Vals	InfoGain	GainRatio
1          3	0.030971064	0.024891536
    2     50	0.060860038	0.027552191
    3     19	0.039562857	0.018676787
    4     37	0.052541227	0.052620750
    5      3	0.074523225	0.047699231
    6     61	0.10604433	0.024471911
    7     20	0.12348668	0.034953203
    8     69	0.097198760	0.043983864
    9      2	0.045752381	0.046816705
   10     64	0.21388759	0.042844587
   11     18	0.66970458	0.18507018
   12     43	1.2780762	0.32537181

-test2-Preparation took 0 seconds, 12 milliseconds and 393 microseconds
Feature Permutation based on GainRatio/Values :
< 9, 5, 11, 1, 12, 7, 4, 3, 10, 8, 2, 6 >
-test2-Phase 2: Learning from Datafile: dimin.train
-test2-Start:          0 @ Fri Jul  3 10:00:23 2026
-test2-Finished:     100 @ Fri Jul  3 10:00:23 2026

Size of InstanceBase = 954 Nodes, (38160 bytes), 26.62 % compression
-test2-Learning took 0 seconds, 0 milliseconds and 409 microseconds
-test2-Phase 2: Appending from Datafile: dimin.train (starting at line 101)
-test2-Start:        101 @ Fri Jul  3 10:00:23 2026
-test2-Learning:     100 @ Fri Jul  3 10:00:23 2026	 added:9
-test2-Learning:     101 @ Fri Jul  3 10:00:23 2026	 added:0
-test2-Learning:     102 @ Fri Jul  3 10:00:23 2026	 added:0
-test2-Learning:     103 @ Fri Jul  3 10:00:23 2026	 added:0
-test2-Learning:     104 @ Fri Jul  3 10:00:23 2026	 added:0
-test2-Learning:     105 @ Fri Jul  3 10:00:23 2026	 added:0
-test2-Learning:     106 @ Fri Jul  3 10:00:23 2026	 added:0
-test2-Learning:     107 @ Fri Jul  3 10:00:23 2026	 added:0
-test2-Learning:     108 @ Fri Jul  3 10:00:23 2026	 added:0
-test2-Learning:     109 @ Fri Jul  3 10:00:23 2026	 added:0
-test2-Learning:     110 @ Fri Jul  3 10:00:23 2026	 added:0
-test2-Learning:     200 @ Fri Jul  3 10:00:23 2026	 added:5
-test2-Learning:    1100 @ Fri Jul  3 10:00:23 2026	 added:65
-test2-Finished:    2898 @ Fri Jul  3 10:00:23 2026
in total added 173 new entries                                      [Note 2]

Size of InstanceBase = 2232 Nodes, (89280 bytes), 32.40 % compression
DB Entropy        : 1.61789286
Number of Classes : 5

Feats	Vals	InfoGain	GainRatio
    1      3	0.03097106	0.02489154
    2     50	0.06086004	0.02755219
    3     19	0.03956286	0.01867679
    4     37	0.05254123	0.05262075
    5      3	0.07452322	0.04769923
    6     61	0.10604433	0.02447191
    7     20	0.12348668	0.03495320
    8     69	0.09719876	0.04398386
    9      2	0.04575238	0.04681670
   10     64	0.21388759	0.04284459
   11     18	0.66970458	0.18507018
   12     43	1.27807625	0.32537181

Examine datafile 'dimin.test' gave the following results:
Number of Features: 12
InputFormat       : C4.5


Starting to test, Testfile: dimin.test
Writing output in:          my_second_test.out
Algorithm     : IB2
Global metric : Overlap
Deviant Feature Metrics:(none)
Weighting     : GainRatio
Feature 1	 : 0.026241147173103
Feature 2	 : 0.030918769841214
Feature 3	 : 0.021445836516602
Feature 4	 : 0.056561885447060
Feature 5	 : 0.048311436541460
Feature 6	 : 0.027043360641622
Feature 7	 : 0.037453180788027
Feature 8	 : 0.044999091421718
Feature 9	 : 0.048992032381874
Feature 10	 : 0.044544230779268
Feature 11	 : 0.185449683494634
Feature 12	 : 0.324719540921155

-test2-Tested:      1 @ Fri Jul  3 10:00:23 2026
-test2-Tested:      2 @ Fri Jul  3 10:00:23 2026
-test2-Tested:      3 @ Fri Jul  3 10:00:23 2026
-test2-Tested:      4 @ Fri Jul  3 10:00:23 2026
-test2-Tested:      5 @ Fri Jul  3 10:00:23 2026
-test2-Tested:      6 @ Fri Jul  3 10:00:23 2026
-test2-Tested:      7 @ Fri Jul  3 10:00:23 2026
-test2-Tested:      8 @ Fri Jul  3 10:00:23 2026
-test2-Tested:      9 @ Fri Jul  3 10:00:23 2026
-test2-Tested:     10 @ Fri Jul  3 10:00:23 2026
-test2-Tested:    100 @ Fri Jul  3 10:00:23 2026
-test2-Ready:     950 @ Fri Jul  3 10:00:23 2026
Seconds taken: 0.1065 (8922.79 p/s)

overall accuracy:        0.940000  (893/950), of which 15 exact matches
There were 43 ties of which 31 (72.09%) were correctly resolved
                                                         [Note 3]
\end{verbatim}
\end{footnotesize}


Notes:
\begin{enumerate}
\item IB2 is bootstrapped with 100 lines, but for the statistics all 2999
 lines are used.
\item As we see here, 173 entries from the input file had a mismatch,
and were therefore entered in the Instancebase.
\item We see that IB2 scores 94.00 \%, compared to 96.21 \% for IGTREE
  in our first example.  For this data, IB2 is not a good
  algorithm. However, it saves a lot of space, and is faster than
  IB1. Yet, IGTREE is both faster and better. Had we used IB1, the
  score would have been 96.84 \%.
\end{enumerate}
\clearpage

\subsection{example 3, {\tt api\_test3.cxx}}

This demonstrates Cross Validation. Let's try the following program:

\begin{footnotesize}
\begin{verbatim}
#include <cstdlib>
#include "timbl/TimblAPI.h"
using Timbl::TimblAPI;

int main(){
  TimblAPI *My_Experiment = new TimblAPI( "-t cross_validate" );
  My_Experiment->Test( "cross_val.test" );
  delete My_Experiment;
  exit(0);
}
\end{verbatim}
\end{footnotesize}

This program creates an experiment, which defaults to IB1 and because of the
special option ``-t cross\_validate'' will start a CrossValidation
experiment.\\
Learn() is not possible now. We must use a special form of Test().

``cross\_val.test'' is a file with the following content:
\begin{footnotesize}
\begin{verbatim}
small_1.train
small_2.train
small_3.train
small_4.train
small_5.train
\end{verbatim}
\end{footnotesize}


All these files contain an equal part of a bigger dataset, and
My\_Experiment will run a CrossValidation test between these files.
Note that output filenames are generated and that you cannot influence
that.

The output of this program is (the folds for {\tt small\_2.train} to
{\tt small\_4.train} are omitted; they are analogous to the first
fold):

\begin{footnotesize}
\begin{verbatim}
Starting Cross validation test on files:
small_1.train
small_2.train
small_3.train
small_4.train
small_5.train
Examine datafile 'small_1.train' gave the following results:
Number of Features: 8
InputFormat       : C4.5


Starting to test, Testfile: small_1.train
Writing output in:          small_1.train.cv
Algorithm     : CV
Global metric : Overlap
Deviant Feature Metrics:(none)
Weighting     : GainRatio

Tested:      1 @ Fri Jul  3 10:00:24 2026
Tested:      2 @ Fri Jul  3 10:00:24 2026
Tested:      3 @ Fri Jul  3 10:00:24 2026
Tested:      4 @ Fri Jul  3 10:00:24 2026
Tested:      5 @ Fri Jul  3 10:00:24 2026
Tested:      6 @ Fri Jul  3 10:00:24 2026
Tested:      7 @ Fri Jul  3 10:00:24 2026
Tested:      8 @ Fri Jul  3 10:00:24 2026
Tested:      9 @ Fri Jul  3 10:00:24 2026
Tested:     10 @ Fri Jul  3 10:00:24 2026
Ready:      10 @ Fri Jul  3 10:00:24 2026
Seconds taken: 0.0026 (3785.01 p/s)

overall accuracy:        0.800000  (8/10)

[... folds 2, 3 and 4: accuracies 0.800000, 0.900000, 0.800000 ...]

Examine datafile 'small_5.train' gave the following results:
Number of Features: 8
InputFormat       : C4.5


Starting to test, Testfile: small_5.train
Writing output in:          small_5.train.cv
Algorithm     : CV
Global metric : Overlap
Deviant Feature Metrics:(none)
Weighting     : GainRatio

Tested:      1 @ Fri Jul  3 10:00:24 2026
Tested:      2 @ Fri Jul  3 10:00:24 2026
Tested:      3 @ Fri Jul  3 10:00:24 2026
Tested:      4 @ Fri Jul  3 10:00:24 2026
Tested:      5 @ Fri Jul  3 10:00:24 2026
Tested:      6 @ Fri Jul  3 10:00:24 2026
Tested:      7 @ Fri Jul  3 10:00:24 2026
Tested:      8 @ Fri Jul  3 10:00:24 2026
Ready:       8 @ Fri Jul  3 10:00:24 2026
Seconds taken: 0.0003 (29962.55 p/s)

overall accuracy:        1.000000  (8/8)
\end{verbatim}
\end{footnotesize}


What has happened here?

\begin{enumerate}
\item TiMBL trained itself with inputfiles small\_2.train through
small\_5.train. (in fact using the {\tt Expand()} API call.
\item Then TiMBL tested small\_1.train against the InstanceBase.
\item Next, small\_2.train is removed from the database (API call {\tt
Remove()} ) and small\_1.train is added.
\item Then small\_2.train is tested against the InstanceBase.
\item And so forth with small\_3.train $\ldots$
\end{enumerate}
\clearpage

\subsection{example 4, {\tt api\_test4.cxx}}

This program demonstrates adding and deleting of the InstanceBase.  It
also proves that weights are (re)calculated correctly each time (which
also explains why this is a time-consuming thing to do). After running
this program, wg.1.wgt should be equal to wg.5.wgt and wg.2.wgt equal to
wg.4.wgt . Important to note is also, that while we do not use a weighting
of X2 or SV here, only the ``simple'' weights are calculated and
stored.

Further, arr.1.arr should be equal to arr.5.arr and arr.2.arr should be equal
to arr.4.arr

First the program:

\begin{footnotesize}
\begin{verbatim}
#include <cstdlib>
#include <iostream>
#include "timbl/TimblAPI.h"
using namespace Timbl;

int main(){
  TimblAPI *My_Experiment = new TimblAPI( "-a IB1 +vDI+DB +mM" ,
                                          "test4" );
  My_Experiment->ShowSettings( std::cout );
  My_Experiment->Learn( "dimin.train" );
  My_Experiment->Test( "dimin.test", "inc1.out" );
  My_Experiment->SaveWeights( "wg.1.wgt" );
  My_Experiment->WriteArrays( "arr.1.arr" );
  My_Experiment->Increment( "=,=,=,=,+,k,e,=,-,r,@,l,T" );
  My_Experiment->Test( "dimin.test", "inc2.out" );
  My_Experiment->SaveWeights( "wg.2.wgt" );
  My_Experiment->WriteArrays( "arr.2.arr" );
  My_Experiment->Increment( "+,zw,A,rt,-,k,O,p,-,n,O,n,E" );
  My_Experiment->Test( "dimin.test", "inc3.out" );
  My_Experiment->SaveWeights( "wg.3.wgt" );
  My_Experiment->WriteArrays( "arr.3.arr" );
  My_Experiment->Decrement( "+,zw,A,rt,-,k,O,p,-,n,O,n,E" );
  My_Experiment->Test( "dimin.test", "inc4.out" );
  My_Experiment->SaveWeights( "wg.4.wgt" );
  My_Experiment->WriteArrays( "arr.4.arr" );
  My_Experiment->Decrement( "=,=,=,=,+,k,e,=,-,r,@,l,T" );
  My_Experiment->Test( "dimin.test", "inc5.out" );
  My_Experiment->SaveWeights( "wg.5.wgt" );
  My_Experiment->WriteArrays( "arr.5.arr" );
  delete My_Experiment;
  exit(1);
}
\end{verbatim}
\end{footnotesize}


This produces the following output (the settings dump and the
near-identical blocks for {\tt inc2.out} to {\tt inc4.out} are
omitted):

\begin{footnotesize}
\begin{verbatim}
Examine datafile 'dimin.train' gave the following results:
Number of Features: 12
InputFormat       : C4.5

-test4-Phase 1: Reading Datafile: dimin.train
-test4-Start:          0 @ Fri Jul  3 11:27:09 2026
-test4-Finished:    2999 @ Fri Jul  3 11:27:09 2026
-test4-Calculating Entropy         Fri Jul  3 11:27:09 2026
-test4-Preparation took 0 seconds, 10 milliseconds and 277 microseconds
Feature Permutation based on GainRatio/Values :
< 9, 5, 11, 1, 12, 7, 4, 3, 10, 8, 2, 6 >
-test4-Phase 2: Building multi index on Datafile: dimin.train
-test4-Start:          0 @ Fri Jul  3 11:27:09 2026
-test4-Finished:    2999 @ Fri Jul  3 11:27:09 2026
-test4-
Phase 3: Learning from Datafile: dimin.train
-test4-Start:          0 @ Fri Jul  3 11:27:09 2026
-test4-Finished:    2999 @ Fri Jul  3 11:27:09 2026

Size of InstanceBase = 19231 Nodes, (769240 bytes), 49.77 % compression
-test4-Learning took 0 seconds, 33 milliseconds and 299 microseconds
Examine datafile 'dimin.test' gave the following results:
Number of Features: 12
InputFormat       : C4.5


Starting to test, Testfile: dimin.test
Writing output in:          inc1.out
Algorithm     : IB1
Global metric : Value Difference, Prestored matrix
Deviant Feature Metrics:(none)
Size of value-matrix[1] = 96 Bytes
Size of value-matrix[2] = 704 Bytes
Size of value-matrix[3] = 704 Bytes
Size of value-matrix[4] = 96 Bytes
Size of value-matrix[5] = 96 Bytes
Size of value-matrix[6] = 1496 Bytes
Size of value-matrix[7] = 1496 Bytes
Size of value-matrix[8] = 336 Bytes
Size of value-matrix[9] = 56 Bytes
Size of value-matrix[10] = 2376 Bytes
Size of value-matrix[11] = 1344 Bytes
Size of value-matrix[12] = 936 Bytes
Total Size of value-matrices 9736 Bytes

Weighting     : GainRatio

-test4-Tested:      1 @ Fri Jul  3 11:27:09 2026
-test4-Tested:      2 @ Fri Jul  3 11:27:09 2026
-test4-Tested:      3 @ Fri Jul  3 11:27:09 2026
-test4-Tested:      4 @ Fri Jul  3 11:27:09 2026
-test4-Tested:      5 @ Fri Jul  3 11:27:09 2026
-test4-Tested:      6 @ Fri Jul  3 11:27:09 2026
-test4-Tested:      7 @ Fri Jul  3 11:27:09 2026
-test4-Tested:      8 @ Fri Jul  3 11:27:09 2026
-test4-Tested:      9 @ Fri Jul  3 11:27:09 2026
-test4-Tested:     10 @ Fri Jul  3 11:27:09 2026
-test4-Tested:    100 @ Fri Jul  3 11:27:09 2026
-test4-Ready:     950 @ Fri Jul  3 11:27:09 2026
Seconds taken: 0.0190 (50113.41 p/s)

overall accuracy:        0.963158  (915/950), of which 62 exact matches
There were 6 ties of which 5 (83.33%) were correctly resolved
-test4-Saving Weights in wg.1.wgt
-test4-Saving Probability Arrays in arr.1.arr

[... the blocks for inc2.out, inc3.out and inc4.out are analogous,
     with an Increment() or Decrement() preceding each of them ...]

Starting to test, Testfile: dimin.test
Writing output in:          inc5.out
Algorithm     : IB1
Global metric : Value Difference, Prestored matrix
Deviant Feature Metrics:(none)
Size of value-matrix[1] = 96 Bytes
Size of value-matrix[2] = 704 Bytes
Size of value-matrix[3] = 704 Bytes
Size of value-matrix[4] = 96 Bytes
Size of value-matrix[5] = 96 Bytes
Size of value-matrix[6] = 1496 Bytes
Size of value-matrix[7] = 1496 Bytes
Size of value-matrix[8] = 336 Bytes
Size of value-matrix[9] = 56 Bytes
Size of value-matrix[10] = 2376 Bytes
Size of value-matrix[11] = 1344 Bytes
Size of value-matrix[12] = 936 Bytes
Total Size of value-matrices 9736 Bytes

Weighting     : GainRatio

-test4-Tested:      1 @ Fri Jul  3 11:27:09 2026
-test4-Tested:      2 @ Fri Jul  3 11:27:09 2026
-test4-Tested:      3 @ Fri Jul  3 11:27:09 2026
-test4-Tested:      4 @ Fri Jul  3 11:27:09 2026
-test4-Tested:      5 @ Fri Jul  3 11:27:09 2026
-test4-Tested:      6 @ Fri Jul  3 11:27:09 2026
-test4-Tested:      7 @ Fri Jul  3 11:27:09 2026
-test4-Tested:      8 @ Fri Jul  3 11:27:09 2026
-test4-Tested:      9 @ Fri Jul  3 11:27:09 2026
-test4-Tested:     10 @ Fri Jul  3 11:27:09 2026
-test4-Tested:    100 @ Fri Jul  3 11:27:09 2026
-test4-Ready:     950 @ Fri Jul  3 11:27:09 2026
Seconds taken: 0.0181 (52448.52 p/s)

overall accuracy:        0.963158  (915/950), of which 62 exact matches
There were 6 ties of which 5 (83.33%) were correctly resolved
-test4-Saving Weights in wg.5.wgt
-test4-Saving Probability Arrays in arr.5.arr
\end{verbatim}
\end{footnotesize}
\clearpage

\subsection{example 5, {\tt api\_test5.cxx}}

This program demonstrates the use of neighborSets to classify and
store results. It also demonstrates some neighborSet basics. Note the
use of {\tt icu::UnicodeString} for the instances.

\begin{footnotesize}
\begin{verbatim}
#include <string>
#include "timbl/TimblAPI.h"

using std::endl;
using std::cout;
using std::string;
using namespace Timbl;

int main(){
  TimblAPI *My_Experiment = new TimblAPI( "-a IB1 +vDI+DB+n +mM +k4 " ,
                                          "test5" );
  My_Experiment->Learn( "dimin.train" );
  {
    icu::UnicodeString line =  "=,=,=,=,+,k,e,=,-,r,@,l,T";
    const neighborSet *neighbours1 = My_Experiment->classifyNS( line );
    if ( neighbours1 ){
      cout << "Classify OK on " << line << endl;
      cout << neighbours1;
    }
    else {
      cout << "Classify failed on " << line << endl;
      neighbours1 = new neighborSet();
    }
    neighborSet neighbours2;
    line = "+,zw,A,rt,-,k,O,p,-,n,O,n,E";
    if ( My_Experiment->classifyNS( line, neighbours2 ) ){
      cout << "Classify OK on " << line << endl;
      cout << neighbours2;
    }
    else {
      cout << "Classify failed on " << line << endl;
    }
    line = "+,z,O,n,-,d,A,xs,-,=,A,rm,P";
    const neighborSet *neighbours3 = My_Experiment->classifyNS( line );
    if ( neighbours3 ){
      cout << "Classify OK on " << line << endl;
      cout << neighbours3;
    }
    else {
      cout << "Classify failed on " << line << endl;
      neighbours3 = new neighborSet();
    }
    neighborSet uit2;
    {
      neighborSet uit;
      uit.setShowDistance(true);
      uit.setShowDistribution(true);
      cout << " before first merge " << endl;
      cout << uit;
      uit.merge( *neighbours1 );
      cout << " after first merge " << endl;
      cout << uit;
      uit.merge( *neighbours3 );
      cout << " after second merge " << endl;
      cout << uit;
      uit.merge( neighbours2 );
      cout << " after third merge " << endl;
      cout << uit;
      uit.truncate( 3 );
      cout << " after truncate " << endl;
      cout << uit;
      cout << " test assignment" << endl;
      uit2 = *neighbours1;
    }
    cout << "assignment result: " << endl;
    cout << uit2;
    {
      cout << " test copy construction" << endl;
      neighborSet uit(uit2);
      cout << "result: " << endl;
      cout << uit;
    }
    cout << "almost done!" << endl;
  }
  delete My_Experiment;
  cout << "done!" << endl;
}
\end{verbatim}
\end{footnotesize}

Its expected output is (without further comment):

\begin{footnotesize}
\begin{verbatim}
Examine datafile 'dimin.train' gave the following results:
Number of Features: 12
InputFormat       : C4.5

-test5-Phase 1: Reading Datafile: dimin.train
-test5-Start:          0 @ Fri Jul  3 10:00:25 2026
-test5-Finished:    2999 @ Fri Jul  3 10:00:25 2026
-test5-Calculating Entropy         Fri Jul  3 10:00:25 2026
-test5-Preparation took 0 seconds, 12 milliseconds and 787 microseconds
Feature Permutation based on GainRatio/Values :
< 9, 5, 11, 1, 12, 7, 4, 3, 10, 8, 2, 6 >
-test5-Phase 2: Building multi index on Datafile: dimin.train
-test5-Start:          0 @ Fri Jul  3 10:00:25 2026
-test5-Finished:    2999 @ Fri Jul  3 10:00:25 2026
-test5-
Phase 3: Learning from Datafile: dimin.train
-test5-Start:          0 @ Fri Jul  3 10:00:25 2026
-test5-Finished:    2999 @ Fri Jul  3 10:00:25 2026

Size of InstanceBase = 19231 Nodes, (769240 bytes), 49.77 % compression
-test5-Learning took 0 seconds, 29 milliseconds and 875 microseconds
Classify OK on =,=,=,=,+,k,e,=,-,r,@,l,T
# k=1	{ T 1.00000 }	0.00000000000000
# k=2	{ T 1.00000 }	0.00318629024734
# k=3	{ T 1.00000 }	0.00341823151183
# k=4	{ T 1.00000 }	0.00374337728446
Classify OK on +,zw,A,rt,-,k,O,p,-,n,O,n,E
# k=1	{ E 1.00000 }	0.00000000000000
# k=2	{ E 1.00000 }	0.05666788032719
# k=3	{ E 1.00000 }	0.06255263661774
# k=4	{ E 1.00000 }	0.06442386036189
Classify OK on +,z,O,n,-,d,A,xs,-,=,A,rm,P
# k=1	{ P 1.00000 }	0.05972983625517
# k=2	{ P 1.00000 }	0.08774076913265
# k=3	{ P 1.00000 }	0.08844278891972
# k=4	{ P 1.00000 }	0.09705864995143
 before first merge
 after first merge
# k=1	{ P 1.00000 }	0.05972983625517
# k=2	{ P 1.00000 }	0.08774076913265
# k=3	{ P 1.00000 }	0.08844278891972
# k=4	{ P 1.00000 }	0.09705864995143
 after second merge
# k=1	{ P 2.00000 }	0.05972983625517
# k=2	{ P 2.00000 }	0.08774076913265
# k=3	{ P 2.00000 }	0.08844278891972
# k=4	{ P 2.00000 }	0.09705864995143
 after third merge
# k=1	{ E 1.00000 }	0.00000000000000
# k=2	{ E 1.00000 }	0.05666788032719
# k=3	{ P 2.00000 }	0.05972983625517
# k=4	{ E 1.00000 }	0.06255263661774
# k=5	{ E 1.00000 }	0.06442386036189
# k=6	{ P 2.00000 }	0.08774076913265
# k=7	{ P 2.00000 }	0.08844278891972
# k=8	{ P 2.00000 }	0.09705864995143
 after truncate
# k=1	{ E 1.00000 }	0.00000000000000
# k=2	{ E 1.00000 }	0.05666788032719
# k=3	{ P 2.00000 }	0.05972983625517
 test assignment
assignment result:
# k=1	{ P 1.00000 }	0.05972983625517
# k=2	{ P 1.00000 }	0.08774076913265
# k=3	{ P 1.00000 }	0.08844278891972
# k=4	{ P 1.00000 }	0.09705864995143
 test copy construction
result:
# k=1	{ P 1.00000 }	0.05972983625517
# k=2	{ P 1.00000 }	0.08774076913265
# k=3	{ P 1.00000 }	0.08844278891972
# k=4	{ P 1.00000 }	0.09705864995143
almost done!
done!
\end{verbatim}
\end{footnotesize}
\clearpage

\subsection{example 6, {\tt api\_test6.cxx}}

This program demonstrates the use of ClassDistributions, TargetValues
and neighborSets for classification. Note the new iterator interface
({\tt *it}, {\tt it->Value()}, {\tt it->Weight()}) and the use of a
normalisation ({\tt -G 0}), which makes the distribution weights
probabilities.

\begin{footnotesize}
\begin{verbatim}
#include <iostream>
#include "timbl/TimblAPI.h"

using std::cout;
using std::endl;
using namespace Timbl;

int main(){
  TimblAPI My_Experiment( "-a IB1 +vDI+DB -G 0 -k3", "test6" );
  My_Experiment.Learn( "dimin.train" );
  const ClassDistribution *vd;
  const TargetValue *tv
    = My_Experiment.Classify( std::string("-,=,O,m,+,h,K,=,-,n,I,N,K"), vd );
  cout << "resulting target: " << tv << endl;
  cout << "resulting Distribution: " << vd << endl;
  ClassDistribution::dist_iterator it=vd->begin();
  while ( it != vd->end() ){
    cout << *it << " OR ";
    cout << it->Value() << " " << it->Weight() << endl;
    ++it;
  }

  cout << "the same with neighborSets" << endl;
  const neighborSet *nb = My_Experiment.classifyNS( "-,=,O,m,+,h,K,=,-,n,I,N,K" );
  WClassDistribution *vd2 = nb->bestDistribution();
  vd2->Normalize();
  cout << "default answer " << vd2 << endl;
  decayStruct *dc = new  expDecay(0.3);
  delete vd2;
  vd2 = nb->bestDistribution( dc );
  delete dc;
  cout << "with exponenial decay, alpha = 0.3 " << vd2 << endl;
  delete vd2;
}
\end{verbatim}
\end{footnotesize}

This is the output produced:

\begin{footnotesize}
\begin{verbatim}
Examine datafile 'dimin.train' gave the following results:
Number of Features: 12
InputFormat       : C4.5

-test6-Phase 1: Reading Datafile: dimin.train
-test6-Start:          0 @ Fri Jul  3 11:27:14 2026
-test6-Finished:    2999 @ Fri Jul  3 11:27:14 2026
-test6-Calculating Entropy         Fri Jul  3 11:27:14 2026
-test6-Preparation took 0 seconds, 11 milliseconds and 245 microseconds
Feature Permutation based on GainRatio/Values :
< 9, 5, 11, 1, 12, 7, 4, 3, 10, 8, 2, 6 >
-test6-Phase 2: Building multi index on Datafile: dimin.train
-test6-Start:          0 @ Fri Jul  3 11:27:14 2026
-test6-Finished:    2999 @ Fri Jul  3 11:27:14 2026
-test6-
Phase 3: Learning from Datafile: dimin.train
-test6-Start:          0 @ Fri Jul  3 11:27:14 2026
-test6-Finished:    2999 @ Fri Jul  3 11:27:14 2026

Size of InstanceBase = 19231 Nodes, (769240 bytes), 49.77 % compression
-test6-Learning took 0 seconds, 29 milliseconds and 168 microseconds
resulting target: K
resulting Distribution: { E 0.125000, K 0.875000 }
E 0.125000 OR E 0.125000
K 0.875000 OR K 0.875000
the same with neighborSets
default answer { E 0.125000, K 0.875000 }
with exponenial decay, alpha = 0.3 { E 0.971556, K 6.69810 }
\end{verbatim}
\end{footnotesize}

\clearpage

\begin{thebibliography}{1}

\bibitem{Daelemans+26}
Walter Daelemans, Jakub Zavrel, Ko van der Sloot, and Antal van den
  Bosch (2026).
\newblock {\em {TiMBL}: Tilburg Memory-Based Learner, version 7.0,
  Reference Guide}.
\newblock Available from \timblmanualurl.

\end{thebibliography}

\end{document}
